\noindent УДК 621.396.96:629.78

\vspace{1em}

\begin{center}
    {\bfseries\large Название статьи\par}
    \vspace{1em}
    {\bfseries И.\,А.~Мартыненко}\par
    {\small НИУ МЭИ, Москва, Россия\par}
    {\small E-mail: ivan.martynenko@proton.me\par}
    {\small ORCID: \href{https://orcid.org/0009-0005-5176-0151}{0009-0005-5176-0151}\par}
\end{center}

\vspace{1em}
\begin{abstract}
\noindent В работе представлена методология сквозного проектирования и тактово-точной верификации приемного тракта первичного поиска сигналов GPS L1 без использования проприетарных САПР. Аппаратный вычислитель алгоритма параллельного поиска (PCPS) разработан на синтезируемом подмножестве SystemVerilog, а тестовое окружение реализовано на базе компилятора Verilator с применением C++. 

Предложенный подход объединяет RTL-модель и программную среду в единый исполняемый бинарный файл, обеспечивая полную автоматизацию сквозного моделирования. Использование Verilator гарантирует абсолютную гибкость верификационного стенда: высокоуровневый C++ позволяет бесшовно подключать сложные алгоритмы пороговой обработки, расчет статистических параметров шума и внешние математические библиотеки напрямую к аппаратному ядру. Данный метод на порядки превосходит классические симуляторы по скорости моделирования, снижает зависимость от закрытых инструментов вендоров и ориентирован на последующую интеграцию в заказные СБИС и платформы RISC-V.
\end{abstract}

\noindent \textbf{Ключевые слова:} GPS L1, первичный поиск, PCPS, критерий Неймана--Пирсона, Verilator, CubeSat, SystemVerilog, HW/SW Co-design, ASIC.

\vspace{1.5em}
\vspace{1.5em}
